% Source: Wald GR, physical PDF pp. 193--195
%         (printed pp. 180--182).
% Coverage: Section 7.4, Methods for Generating Solutions;
%           equations (7.4.1)--(7.4.14).
% Figures/tables/footnotes: footnote 5; no figures or tables.
% Status: source-reviewed and compile-clean.
% Uncertainties: none; a printed cross-reference error is corrected below.

\subsection{Methods for Generating Solutions}\label{sec:7.4}

It often happens when one is trying to solve an equation that an algorithm will exist
for constructing new solutions from a given solution. For example, for Laplace's
equation in ordinary electrostatics one can construct new solutions from a given
solution by the “method of inversion” (see, e.g., Jackson 1962).\footnote{The
method of inversion works because sphere inversion maps the Euclidean metric into
a multiple of itself (i.e., it is a “conformal isometry”) and Laplace's equation has
simple conformal transformation properties (see appendix D).} Although such
prescriptions often can be applied readily, they usually suffer from the serious
drawback that the solutions they generate may be of no physical interest, or at least,
not applicable to the problem that one desires to solve.

For the vacuum Einstein equation, \(R_{ab}=0\), it turns out that if one is given a
solution with a Killing vector field \(\xi^a\) there exists an algorithm for
constructing a one-parameter family of solutions, for which \(\xi^a\) remains a
Killing field. Algorithms of this sort have been given by Ehlers (1957) and others.
We shall present the algorithm in the more general form given by Geroch (1971). To
define it, we need to use the fact that as proven in section 7.1 (see
equation~\eqref{eq:7.1.15}) in a vacuum spacetime, the \emph{twist},
\(\omega_a\), of a Killing field \(\xi^a\), defined by
\begin{equation}
  \omega_a=\epsilon_{abcd}\xi^b\nabla^c\xi^d ,
  \tag{7.4.1}\label{eq:7.4.1}
\end{equation}
satisfies
\begin{equation}
  \nabla_{[a}\omega_{b]}=0 .
  \tag{7.4.2}\label{eq:7.4.2}
\end{equation}
Consequently, as remarked at the end of section B.1 of appendix B, locally there
exists a scalar function, \(\omega\), such that
\begin{equation}
  \omega_a=\nabla_a\omega
  \tag{7.4.3}\label{eq:7.4.3}
\end{equation}
(\(\omega\) is called the \emph{scalar twist} of \(\xi^a\)). Furthermore, using
equation~\eqref{eq:B.2.13}, we find
\begin{equation}
  \epsilon^{abef}\nabla_e
    \left[\epsilon_{abcd}\nabla^c\xi^d\right]
    =-4\nabla_e\nabla^e\xi^f .
  \tag{7.4.4}\label{eq:7.4.4}
\end{equation}
Thus, using equation~\eqref{eq:C.3.9} and the vacuum Einstein equation
\(R_{ab}=0\), we find
\begin{equation}
  \nabla_{[e}
    \left\{\epsilon_{ab]cd}\nabla^c\xi^d\right\}=0 ,
  \tag{7.4.5}\label{eq:7.4.5}
\end{equation}
which implies the (local) existence of a one-form field \(\alpha_b\) such that
\begin{equation}
  \nabla_{[a}\alpha_{b]}
    =\frac12\epsilon_{abcd}\nabla^c\xi^d .
  \tag{7.4.6}\label{eq:7.4.6}
\end{equation}
By adding a gradient to \(\alpha_b\), we can adjust it so that
\begin{equation}
  \xi^a\alpha_a=\omega .
  \tag{7.4.7}\label{eq:7.4.7}
\end{equation}
Finally, similar calculations (problem 5) show that
\begin{equation}
  \nabla_{[e}
    \left\{
      2\lambda\nabla_a\xi_{b]}
      +\omega\epsilon_{ab]cd}\nabla^c\xi^d
    \right\}=0 ,
  \tag{7.4.8}\label{eq:7.4.8}
\end{equation}
where
\begin{equation}
  \lambda=\xi^c\xi_c ,
  \tag{7.4.9}\label{eq:7.4.9}
\end{equation}
and thus, locally, there exists a one-form \(\beta_b\) such that
\begin{equation}
  \nabla_{[a}\beta_{b]}
    =2\lambda\nabla_a\xi_b
      +\omega\epsilon_{abcd}\nabla^c\xi^d ,
  \tag{7.4.10}\label{eq:7.4.10}
\end{equation}
where \(\beta_b\) can be adjusted to satisfy
\begin{equation}
  \xi^a\beta_a=\omega^2+\lambda^2-1 .
  \tag{7.4.11}\label{eq:7.4.11}
\end{equation}

The algorithm for generating new solutions from a given solution now may be stated
as follows: Let \(g_{ab}\) be a solution of the vacuum Einstein equation \(R_{ab}=0\)
with Killing field \(\xi^a\). Define \(\omega\), \(\alpha_a\), and \(\beta_a\) as
above. For each \(\theta\in[0,\pi]\), define \(\widetilde g_{ab}(\theta)\) by
\begin{equation}
  \widetilde g_{ab}(\theta)
    =\sigma g_{ab}
      +2\sin\theta\,\xi_{(a}\gamma_{b)}
      +\frac{\lambda\sin^2\theta}{\sigma}\gamma_a\gamma_b ,
  \tag{7.4.12}\label{eq:7.4.12}
\end{equation}
where \(\sigma(\theta)\) and \(\gamma_a(\theta)\) are defined by
\begin{align}
  \sigma&=(\cos\theta-\omega\sin\theta)^2
    +\lambda^2\sin^2\theta ,
  \tag{7.4.13}\label{eq:7.4.13}\\
  \gamma_a&=2\alpha_a\cos\theta-\beta_a\sin\theta .
  \tag{7.4.14}\label{eq:7.4.14}
\end{align}
\emph{Then for all \(\theta\), \(\widetilde g_{ab}(\theta)\) is a solution of the
vacuum Einstein equation.} Furthermore, \(\xi^a\) remains a Killing field of
\(\widetilde g_{ab}(\theta)\). Note that for \(\theta=\pi\),
\(\widetilde g_{ab}\) reduces to \(g_{ab}\); however, otherwise
\(\widetilde g_{ab}(\theta)\) is, in general, a distinct solution. Note also that
even if \(g_{ab}\) is nonsingular, \(\widetilde g_{ab}\) may develop singularities
on account of the bad behavior of \(\omega\), \(\alpha_a\), or \(\beta_a\)
resulting from the possibility that one cannot define them globally.

The reader may verify by direct computation that equation~\eqref{eq:7.4.12} indeed
does define a solution of the vacuum Einstein equation. However, for a derivation of
equation~\eqref{eq:7.4.12} which gives more insight into why this algorithm works,
we refer the reader to Geroch (1971).

Unfortunately, the algorithm~\eqref{eq:7.4.12} does not reproduce most physically
interesting features of the original metric. In particular, \(\widetilde g_{ab}\) will
not, in general, be asymptotically flat even if \(g_{ab}\) satisfies this property.
Furthermore, equation~\eqref{eq:7.4.12} produces only a one-parameter family of
solutions. If one reapplies the algorithm to \(\widetilde g_{ab}(\theta)\), one does
not generate any new solutions.

However, if \(g_{ab}\) has two commuting Killing vector fields, \(\xi^a\) and
\(\psi^a\) (as is the case in the stationary axisymmetric spacetimes considered in
section 7.1), it can be shown (Geroch 1972a) that \(\psi^a\) also remains a Killing
vector field of the metric \(\widetilde g_{ab}(\theta)\) generated by the
algorithm~\eqref{eq:7.4.12} using the Killing field \(\xi^a\). Hence, we may apply
the algorithm to \(\widetilde g_{ab}(\theta)\) using the Killing vector \(\psi^a\)
to generate a two-parameter family of solutions
\(\widetilde g_{ab}(\theta,\theta')\) from our original solution \(g_{ab}\). We
then may reapply the transformation~\eqref{eq:7.4.12} to
\(\widetilde g_{ab}(\theta,\theta')\). Geroch (1972a) has shown that this
procedure does not, in general, reproduce previous solutions, and thus a
three-parameter family of solutions is generated. Indeed, an infinite-dimensional
group of transformations is generated by repeated application of the transformation
% SOURCE ERRATUM: printed ``(7.2.12)''; corrected to (7.4.12).
\eqref{eq:7.4.12}. Recently, it has been proven (Hauser and Ernst 1981) that---as conjectured
by Geroch (1972a)---all asymptotically flat, stationary, axisymmetric vacuum
solutions of Einstein's equation can be generated from Minkowski spacetime by an
element of this group. Furthermore, it is known how to generate such asymptotically
flat solutions with desired values for all multipole moments (Hoenselaers, Kinnersley,
and Xanthopoulos 1979; Xanthopoulos 1981). Unfortunately, the algebraic
computations required to obtain these solutions explicitly remain formidable, so in
fact very few solutions are presently known in explicit form. Nevertheless, this
method for generating solutions is unquestionably one of the most important
techniques for solving Einstein's equation.
